\documentclass[11pt]{article}

\usepackage[a4paper,margin=1in]{geometry}
\usepackage{amsmath,amssymb,amsthm,mathtools}
\usepackage{enumitem}
\usepackage{hyperref}
\usepackage{xcolor}

\newtheorem{definition}{Definition}
\newtheorem{problem}{Problem}
\newtheorem{hypothesis}{Hypothesis}
\newtheorem{proposition}{Proposition}
\newtheorem{experiment}{Experiment}
\newtheorem{remark}{Remark}

\newcommand{\Vocab}{\mathcal V}
\newcommand{\Mask}{\mathtt{MASK}}
\newcommand{\R}{\mathbb R}
\newcommand{\E}{\mathbb E}
\newcommand{\TT}{\mathrm{TT}}
\newcommand{\TN}{\mathrm{TN}}
\newcommand{\KL}{\mathrm{KL}}
\newcommand{\Ent}{\mathrm{Ent}}
\newcommand{\Sat}{\mathrm{Sat}}
\newcommand{\Law}{\mathcal L}
\newcommand{\cC}{\mathcal C}
\newcommand{\cM}{\mathcal M}
\newcommand{\cG}{\mathcal G}
\newcommand{\cF}{\mathcal F}
\newcommand{\cX}{\mathcal X}
\newcommand{\cE}{\mathcal E}
\newcommand{\dd}{\,\mathrm d}

\title{Diffusion Language Models and Logical Tensor Networks:\\Research Directions and First Experiments}
\author{}
\date{\today}

\begin{document}
\maketitle

\begin{abstract}
Masked diffusion language models generate sequences by iterative denoising rather than purely left-to-right autoregressive factorization. This makes them natural candidates for globally constrained generation, infilling, repair, and recursive refinement. Tensor methods and logical tensor networks provide complementary tools: low-rank representation of high-dimensional discrete states, tensorized factor-graph inference, differentiable logic, and compact policies over combinatorial masking decisions. This note identifies promising research directions at the intersection of diffusion LLMs, tensor trains, tensor-network reasoning, and logical tensor networks. It also proposes initial experiments using \texttt{tinyTT} and \texttt{tnreason}.
\end{abstract}

\section{Motivation}

A masked diffusion language model (MDLM) maintains a partially specified sequence
\[
    x^{(k)} \in (\Vocab \cup \{\Mask\})^n
\]
and iteratively denoises it until a fully specified sequence is obtained. Unlike an autoregressive model, the MDLM is not forced to decide tokens in left-to-right order. It can refine arbitrary positions, use bidirectional context, and remask or repair inconsistent spans.

This suggests a connection to adaptive numerical solvers and tensor-network inference. The object to be generated is a high-dimensional discrete configuration
\[
    x = (x_1,\ldots,x_n) \in \Vocab^n,
\]
and denoising resembles approximate inference in a complex conditional distribution
\[
    p_\theta(x \mid c),
\]
where \(c\) denotes the prompt or conditioning context. Logical constraints, type constraints, syntax constraints, theorem dependencies, proof obligations, or program tests may be represented as additional factors. Tensor networks provide compressed representations of such high-dimensional factors and of the corresponding approximate marginals.

The central research hypothesis is the following.

\begin{hypothesis}[Tensor-structured recursive diffusion]
Masked diffusion LMs can be interpreted and improved as recursive low-rank inference procedures over partially specified symbolic tensor states. Tensor networks and logical tensor networks can be used to represent constraint structure, guide denoising order, compress mutable memory, and define verifier-guided correction policies.
\end{hypothesis}

\section{Background}

\subsection{Masked diffusion language models}

Let \(x_0 \in \Vocab^n\) denote a clean sequence. A corruption process samples a masked sequence \(x_t\) at noise level \(t\), for example by independently replacing positions by \(\Mask\). The denoising model predicts the clean token distribution at masked positions:
\[
    p_\theta(x_{0,i} \mid x_t,c,t), \qquad i \in M_t,
\]
where \(M_t = \{ i : x_{t,i} = \Mask \}\).

A standard training loss is a weighted masked-token objective
\[
    \mathcal L_{\mathrm{denoise}}(\theta)
    =
    \E_{x_0,t,x_t}
    \left[
        \sum_{i\in M_t} w(t)
        \bigl(-\log p_\theta(x_{0,i}\mid x_t,c,t)\bigr)
    \right].
\]
At inference time, one starts from a heavily masked sequence and repeatedly unmaskes positions according to a schedule. The schedule may be confidence-based, entropy-based, random, blockwise, or learned.

The important degree of freedom is the mask policy
\[
    \pi_\eta(M_{k+1} \mid x^{(k)},c),
\]
which decides which locations should be filled, preserved, or remasked.

\subsection{Tensor-network view of discrete reasoning}

A symbolic reasoning problem can often be written as inference over discrete variables
\[
    z = (z_1,\ldots,z_d), \qquad z_j \in \mathcal A_j.
\]
Logical clauses, compatibility constraints, grammar rules, or probabilistic relations define factors
\[
    \psi_a(z_{\partial a}) \ge 0,
\]
and a joint unnormalized score
\[
    \Psi(z) = \prod_{a\in \cF} \psi_a(z_{\partial a}).
\]
The corresponding partition function and marginals are tensor contractions:
\[
    Z = \sum_{z_1,\ldots,z_d} \prod_{a\in \cF} \psi_a(z_{\partial a}).
\]
If the factor graph has exploitable structure, the contraction may be approximated or represented in a tensor-train, tree tensor network, hierarchical Tucker, or general factor-graph form.

\subsection{Logical tensor networks}

Logical Tensor Networks represent predicates, constants, functions, and first-order formulas in differentiable real-valued semantics. A formula \(\varphi\) receives a soft truth value
\[
    \Sat_\varphi(u) \in [0,1],
\]
and a knowledge base \(\Phi = \{\varphi_1,\ldots,\varphi_m\}\) receives an aggregate satisfaction score
\[
    \Sat_\Phi(u)
    =
    \operatorname{Agg}\bigl(
        \Sat_{\varphi_1}(u),\ldots,\Sat_{\varphi_m}(u)
    \bigr).
\]
For diffusion LMs, the satisfaction score can be used as an auxiliary loss, an inference-time energy, a verifier signal, or a remasking rule.

\section{Research direction I: Verifier-guided recursive masked diffusion}

\subsection{Core idea}

Use a base masked diffusion LM as a proposal mechanism and combine it with an external symbolic or tensor-network verifier. The loop is:
\[
\begin{aligned}
    x^{(k+1/2)} &\sim K_\theta(\cdot \mid x^{(k)},c,M_k),\\
    r^{(k+1/2)} &= \operatorname{Verifier}(x^{(k+1/2)},c),\\
    M_{k+1} &= \operatorname{MaskPolicy}(r^{(k+1/2)},x^{(k+1/2)},c),\\
    x^{(k+1)} &= \operatorname{Mask}(x^{(k+1/2)},M_{k+1}).
\end{aligned}
\]
Here \(r\) is a residual-like object: syntax error locations, failed tests, violated clauses, inconsistent proof dependencies, unsatisfied predicates, or low-confidence logical roles.

\subsection{Why this is promising}

This uses the native strength of MDLMs: arbitrary-position repair. An autoregressive LM can append corrections, but a diffusion LM can directly remask and regenerate the defective region. This is particularly attractive for:
\begin{itemize}[leftmargin=2em]
    \item program repair and code infilling;
    \item theorem/proof sketch repair;
    \item formal-language tasks;
    \item Sudoku, Countdown, SAT, and CSPs;
    \item long structured documents with global consistency constraints;
    \item scientific writing where notation, assumptions, constants, and theorem dependencies must remain consistent.
\end{itemize}

\subsection{Idealized convergence statement}

Let \(V:\Vocab^n\to [0,\infty)\) measure constraint violation. Suppose that for the chosen mask \(M_k\), the diffusion repair kernel satisfies
\[
    \E\bigl[V(x^{(k+1/2)}) \mid x^{(k)},M_k\bigr]
    \le
    (1-\rho(M_k;x^{(k)})) V(x^{(k)}) + \varepsilon_\theta.
\]
Assume the mask policy is residual-effective:
\[
    \rho(M_k;x^{(k)}) \ge \rho_0 > 0
\]
whenever \(V(x^{(k)})>\tau\). Then
\[
    \E[V(x^{(K)})]
    \le
    (1-\rho_0)^K V(x^{(0)})
    +
    \frac{\varepsilon_\theta}{\rho_0}.
\]
This theorem is deliberately stylized. Its role is to identify the core mechanism: verifier-guided remasking should reduce a residual-like violation functional up to model error.

\subsection{Initial experiment}

Choose a finite symbolic task with a cheap verifier, for example Sudoku, Boolean formulas, simple typed expressions, or bracketed arithmetic expressions.

\paragraph{Data.}
Generate clean solutions \(x_0\), corrupt them by masking spans, and train or emulate a denoising model. For a first test, the denoiser can be simple:
\[
    p_\theta(x_i \mid x_{-i}) \propto \exp(s_\theta(i,x_i,x_{\mathrm{obs}})).
\]
It need not yet be a large LM.

\paragraph{Verifier.}
Implement a verifier producing local residuals:
\[
    r_j(x) \in \{0,1\}, \qquad j=1,\ldots,J,
\]
where \(r_j(x)=1\) means constraint \(j\) is violated.

\paragraph{Masking policies.}
Compare:
\begin{enumerate}[leftmargin=2em]
    \item random remasking;
    \item confidence/entropy remasking;
    \item verifier-local remasking;
    \item tensor-network marginal remasking;
    \item learned TT mask policy.
\end{enumerate}

\paragraph{Metrics.}
Measure success rate, number of denoising steps, violation decay \(V(x^{(k)})\), and wall-clock cost.

\section{Research direction II: Tensor-network mask policies}

\subsection{Problem}

The mask policy is a combinatorial object. For a sequence of length \(n\), a mask is
\[
    m = (m_1,\ldots,m_n)\in\{0,1\}^n.
\]
A general policy over masks has \(2^n\) entries. Tensor trains provide a compact representation:
\[
    \pi_\eta(m)
    \propto
    \sum_{\alpha_0,\ldots,\alpha_n}
    G_1(\alpha_0,m_1,\alpha_1)
    G_2(\alpha_1,m_2,\alpha_2)
    \cdots
    G_n(\alpha_{n-1},m_n,\alpha_n),
\]
with boundary ranks \(r_0=r_n=1\).

\subsection{Conditioned mask policy}

The policy should depend on features \(f_i\) at each position:
\[
    f_i = \bigl(
        \text{entropy}_i,\,
        \text{confidence}_i,\,
        \text{verifier flags}_i,\,
        \text{syntax role}_i,\,
        \text{step }k
    \bigr).
\]
A simple parameterization is
\[
    \pi_\eta(m\mid f)
    \propto
    \operatorname{TT}_\eta(m_1,\ldots,m_n;f_1,\ldots,f_n).
\]
A more practical first version is an energy model
\[
    E_\eta(m;f)
    =
    \sum_i e_i(m_i;f_i)
    +
    \operatorname{TT}_\eta(m_1,\ldots,m_n),
\]
and
\[
    \pi_\eta(m\mid f)
    \propto
    \exp(-E_\eta(m;f)).
\]

\subsection{Why this is interesting}

A good denoising order is not independent over positions. In reasoning text, one may want:
\begin{itemize}[leftmargin=2em]
    \item premises before derived steps;
    \item definitions before theorem statements;
    \item function signatures before function bodies;
    \item variable declarations before uses;
    \item logical connectives before conclusions;
    \item proof skeleton before local estimates.
\end{itemize}
This is a structured dependency problem. A TT or tree tensor policy can encode long-range dependencies compactly.

\subsection{Initial tinyTT test}

Use \texttt{tinyTT} to represent a binary mask policy for small \(n\), say \(n=16,32,64\).

\paragraph{Synthetic target.}
Define an oracle policy for a formal-language task. Example:
\[
    m_i=1 \quad \text{is preferred if position } i
    \text{ belongs to the next unresolved syntactic block.}
\]
Fit a TT policy to imitate the oracle.

\paragraph{Training loss.}
Given examples \((f^{(s)},m_\star^{(s)})\), minimize
\[
    \mathcal L(\eta)
    =
    -\sum_s \log \pi_\eta(m_\star^{(s)}\mid f^{(s)}).
\]
For a first prototype, avoid stochastic sampling and fit the TT tensor by least squares or cross-entropy over a restricted candidate mask set.

\paragraph{Evaluation.}
Compare the learned TT policy against local independent policies:
\[
    \pi_{\mathrm{local}}(m\mid f)
    =
    \prod_{i=1}^n \pi_i(m_i\mid f_i).
\]
The key question is whether the TT policy captures nonlocal denoising order constraints.

\section{Research direction III: Tensor-network representation of logical constraints}

\subsection{Core idea}

Use \texttt{tnreason} or a similar tensor-network reasoning framework to compile logical constraints into tensor networks. The resulting network defines a soft or hard satisfaction score.

For discrete variables \(z_1,\ldots,z_d\), clauses define factors
\[
    \psi_a(z_{\partial a}) \in \{0,1\}
    \quad\text{or}\quad
    \psi_a(z_{\partial a}) \in [0,1].
\]
The full satisfaction tensor is
\[
    \Psi(z_1,\ldots,z_d)
    =
    \prod_{a\in\cF}
    \psi_a(z_{\partial a}).
\]
Marginals identify promising repairs:
\[
    q_i(v)
    =
    \frac{1}{Z}
    \sum_{z_{-i}}
    \mathbf 1_{\{z_i=v\}}
    \Psi(z).
\]

\subsection{Connection to diffusion denoising}

At a diffusion step, some variables are observed and others are masked. Conditioning on observed tokens gives
\[
    \Psi(z_M \mid z_{\Omega})
    =
    \prod_{a\in\cF}
    \psi_a(z_{\partial a})
    \prod_{i\in\Omega}
    \mathbf 1_{\{z_i=x_i\}}.
\]
Tensor-network contraction yields approximate marginals over masked positions. These marginals can be used to:
\begin{enumerate}[leftmargin=2em]
    \item bias the DLM logits;
    \item select which positions to unmask;
    \item select which positions to remask;
    \item detect contradictions;
    \item produce a logic residual.
\end{enumerate}

\subsection{Logit correction}

Let \(p_\theta(v\mid x^{(k)},c,i)\) be the DLM token distribution at masked position \(i\). Let \(q_i(v)\) be the tensor-network logical marginal. Combine them by product-of-experts:
\[
    \tilde p(v\mid x^{(k)},c,i)
    \propto
    p_\theta(v\mid x^{(k)},c,i)^{\beta}
    q_i(v)^{1-\beta}.
\]
Equivalently,
\[
    \log \tilde p(v)
    =
    \beta \log p_\theta(v)
    +
    (1-\beta)\log q_i(v)
    -
    \log Z_i.
\]
This is a clean and testable mechanism for logic-guided denoising.

\subsection{Initial tnreason test}

Start without a real LLM. Use \texttt{tnreason} to build a small logical tensor network for a symbolic task.

\paragraph{Candidate tasks.}
\begin{itemize}[leftmargin=2em]
    \item Sudoku row/column/box constraints.
    \item Boolean formula completion.
    \item Typed expression completion.
    \item Small proof-state completion.
    \item Parentheses and operator-precedence completion.
\end{itemize}

\paragraph{Protocol.}
\begin{enumerate}[leftmargin=2em]
    \item Generate complete valid assignments.
    \item Mask a subset of variables.
    \item Use tensor-network contraction to compute marginals over masked variables.
    \item Fill high-confidence variables.
    \item Recompute marginals and repeat.
\end{enumerate}

This gives a pure tensor-network analogue of masked diffusion. Then add a learned noisy denoiser and test whether tensor-network marginals improve repair.

\section{Research direction IV: Logical Tensor Network reward or auxiliary loss}

\subsection{Soft satisfaction loss}

Suppose a DLM predicts soft token distributions
\[
    P_i \in \Delta^{|\Vocab|-1}.
\]
Map them to embeddings
\[
    e_i = E^\top P_i.
\]
A differentiable logical tensor network computes a satisfaction score
\[
    \Sat_\Phi(e_1,\ldots,e_n)\in[0,1].
\]
Use the auxiliary objective
\[
    \mathcal L(\theta)
    =
    \mathcal L_{\mathrm{denoise}}(\theta)
    +
    \lambda
    \E\left[
        1-\Sat_\Phi(e_1,\ldots,e_n)
    \right].
\]

\subsection{Practical warning}

Fully differentiable first-order logic over raw BPE tokens is likely brittle. A better route is to apply logical constraints to extracted structured variables:
\[
    h = \operatorname{Extractor}_\psi(P_1,\ldots,P_n),
\]
where \(h\) may contain:
\begin{itemize}[leftmargin=2em]
    \item entity identities;
    \item predicate arguments;
    \item syntax roles;
    \item proof-step labels;
    \item type annotations;
    \item dependency edges.
\end{itemize}
Then evaluate
\[
    \Sat_\Phi(h).
\]

\subsection{Initial LTN-style test}

Train a small denoiser on a symbolic task with known predicates. Example: family-relations reasoning.

Variables:
\[
    \operatorname{Parent}(a,b), \quad
    \operatorname{Ancestor}(a,b).
\]
Axioms:
\[
    \forall a,b:\operatorname{Parent}(a,b)\Rightarrow \operatorname{Ancestor}(a,b),
\]
\[
    \forall a,b,c:
    \operatorname{Parent}(a,b)\wedge \operatorname{Ancestor}(b,c)
    \Rightarrow
    \operatorname{Ancestor}(a,c).
\]
Use an LTN loss to guide the denoiser. Compare:
\begin{enumerate}[leftmargin=2em]
    \item denoising-only training;
    \item denoising plus LTN satisfaction loss;
    \item inference-time tensor-network correction;
    \item combined training and inference correction.
\end{enumerate}

\section{Research direction V: Tensorizing the denoising trajectory}

\subsection{Trajectory tensor}

For a diffusion LM, collect logits across denoising steps:
\[
    \mathcal L(k,i,v)
    =
    \text{logit for token }v
    \text{ at position }i
    \text{ during step }k.
\]
This is a third-order tensor:
\[
    \mathcal L \in \R^{K\times n\times |\Vocab|}.
\]
For hidden states,
\[
    \mathcal H(k,i,j)
    =
    \text{hidden coordinate }j
    \text{ at position }i
    \text{ during step }k.
\]
One can test whether \(\mathcal L\) or \(\mathcal H\) admits low TT, Tucker, CP, or tensor-ring rank.

\subsection{Interpretation}

Low rank in \((k,i,v)\) would suggest that the denoising trajectory has separable structure:
\[
    \mathcal L(k,i,v)
    \approx
    \sum_{\alpha,\beta}
    A(k,\alpha)B(\alpha,i,\beta)C(\beta,v).
\]
This would support the view that diffusion reasoning proceeds through a small number of global modes: outline formation, role assignment, local completion, consistency repair, and final polishing.

\subsection{Initial test}

Use any accessible masked diffusion model or a small trained toy model. Save logits for \(K\) denoising steps. For vocabulary-heavy cases, restrict to:
\begin{itemize}[leftmargin=2em]
    \item top-\(q\) candidate tokens per position;
    \item a symbolic vocabulary;
    \item hidden-state tensors instead of logit tensors.
\end{itemize}

Then run TT-SVD or Tucker decomposition and measure rank/error curves:
\[
    \epsilon(r)
    =
    \frac{
        \|\mathcal L-\mathcal L_r\|_F
    }{
        \|\mathcal L\|_F
    }.
\]
Compare successful and failed reasoning trajectories. The hypothesis is that successful trajectories exhibit cleaner low-rank refinement structure.

\section{Research direction VI: Recursive tensor-diffusion workspaces}

\subsection{Multilevel state}

A recursive LM operates on a mutable workspace rather than a flat sequence. For a long scientific or mathematical document, define levels:
\[
    \ell=0: \text{global problem graph},
\]
\[
    \ell=1: \text{section/lemma/function graph},
\]
\[
    \ell=2: \text{paragraph/proof/code block},
\]
\[
    \ell=3: \text{tokens}.
\]
A tensor tree can represent dependencies between levels:
\[
    \mathcal M
    =
    \operatorname{TTN}
    \bigl(
        \mathcal M_0,\mathcal M_1,\mathcal M_2,\mathcal M_3
    \bigr).
\]

\subsection{Recursive refinement}

At each level:
\[
    x_\ell^{(k+1)}
    =
    \operatorname{Denoise}_\theta
    \bigl(
        x_\ell^{(k)},
        \operatorname{Context}(x_{<\ell}^{(k)},x_{>\ell}^{(k)}),
        r_\ell^{(k)}
    \bigr).
\]
The residual \(r_\ell\) comes from logical, syntactic, or semantic checks. The tensor memory passes coarse information down and aggregates fine information up.

\subsection{Connection to multilevel numerical solvers}

This resembles multigrid or adaptive low-rank correction:
\[
    \text{coarse solve}
    \rightarrow
    \text{prolongation}
    \rightarrow
    \text{local smoothing}
    \rightarrow
    \text{residual correction}.
\]
The analogy is useful because it suggests:
\begin{itemize}[leftmargin=2em]
    \item coarse-to-fine denoising schedules;
    \item residual-based remasking;
    \item rank-adaptive memory compression;
    \item localized correction of inconsistent spans;
    \item contraction estimates for recursive refinement.
\end{itemize}

\section{Recommended first experiments}

\subsection{Experiment A: Pure tensor-network masked inference with tnreason}

\begin{experiment}[Tensor-network denoising without a neural model]
Use \texttt{tnreason} to implement masked completion in a finite logical domain.
\end{experiment}

\paragraph{Goal.}
Show that tensor-network contraction can act as a symbolic denoiser.

\paragraph{Task.}
Start with one of:
\begin{itemize}[leftmargin=2em]
    \item Boolean formula completion;
    \item small Sudoku, e.g. \(4\times 4\);
    \item typed expression completion;
    \item family-relation closure;
    \item finite-state grammar completion.
\end{itemize}

\paragraph{Algorithm.}
\begin{enumerate}[leftmargin=2em]
    \item Compile constraints into a tensor network.
    \item Condition on observed variables.
    \item Contract to obtain marginals over masked variables.
    \item Fill variables with confidence above threshold \(\tau\).
    \item Repeat until completion or contradiction.
\end{enumerate}

\paragraph{Expected result.}
This creates a baseline ``logical diffusion'' process:
\[
    \Mask^n \to x^{(1)} \to x^{(2)} \to \cdots \to x.
\]
It is not an LLM, but it isolates the inference mechanism.

\subsection{Experiment B: Noisy denoiser plus tensor-network correction}

\begin{experiment}[Product-of-experts symbolic denoising]
Combine a noisy learned or heuristic denoiser with tensor-network marginals.
\end{experiment}

Let \(p_\theta(v\mid i,x_{\Omega})\) be a noisy local denoiser and \(q_i(v)\) the tensor-network marginal. Use
\[
    \tilde p_i(v)
    \propto
    p_\theta(v\mid i,x_{\Omega})^\beta q_i(v)^{1-\beta}.
\]
Vary \(\beta\in[0,1]\). This tests how much logical tensor information improves noisy denoising.

\paragraph{Metrics.}
\[
    \text{accuracy},\qquad
    \text{constraint violation},\qquad
    \text{number of steps},\qquad
    \text{contraction cost}.
\]

\subsection{Experiment C: tinyTT mask policy}

\begin{experiment}[TT policy for adaptive remasking]
Use \texttt{tinyTT} to learn or approximate a structured mask policy.
\end{experiment}

\paragraph{Setup.}
For length \(n=16\) or \(n=32\), define oracle masks from a verifier:
\[
    m_\star(x) = \operatorname{locations\_to\_repair}(x).
\]
Fit a TT representation of the mask score:
\[
    S_\eta(m;x) \approx \text{quality of choosing mask }m.
\]

\paragraph{Simplified implementation.}
Avoid full normalization over \(2^n\). Generate a candidate set
\[
    \mathcal C(x) = \{m^{(1)},\ldots,m^{(B)}\}
\]
and train a ranking loss
\[
    \mathcal L(\eta)
    =
    -\log
    \frac{\exp(S_\eta(m_\star;x))}
    {
        \sum_{m\in\mathcal C(x)}
        \exp(S_\eta(m;x))
    }.
\]

\paragraph{Comparison.}
Compare with entropy thresholding and local logistic mask selection.

\subsection{Experiment D: Logic-role denoising order}

\begin{experiment}[Dependency-ordered unmasking]
Classify masked positions by logical role and unmask in dependency order.
\end{experiment}

For synthetic reasoning traces, label positions as:
\[
    \text{premise},\quad
    \text{definition},\quad
    \text{connective},\quad
    \text{derived step},\quad
    \text{conclusion},\quad
    \text{filler}.
\]
Use the schedule:
\[
    \text{premises}
    \to
    \text{definitions}
    \to
    \text{connectives}
    \to
    \text{derived steps}
    \to
    \text{conclusions}
    \to
    \text{filler}.
\]
This is a minimal version of logic-guided diffusion. The tensor-network extension is to infer logical roles and dependencies from a factor graph rather than using fixed labels.

\subsection{Experiment E: Denoising-trajectory tensor rank}

\begin{experiment}[Low-rank structure of refinement trajectories]
Collect denoising trajectories and study tensor rank.
\end{experiment}

For each sample, save logits or hidden states over diffusion steps:
\[
    \mathcal L \in \R^{K\times n\times q}.
\]
Apply TT-SVD or Tucker decomposition. Compare rank/error curves for:
\begin{itemize}[leftmargin=2em]
    \item successful vs failed completions;
    \item random vs logic-guided schedules;
    \item easy vs hard instances;
    \item short vs long reasoning chains.
\end{itemize}

\section{Implementation plan}

\subsection{Milestone 1: tnreason symbolic denoising}

\paragraph{Objective.}
Implement a pure tensor-network masked inference loop.

\paragraph{Steps.}
\begin{enumerate}[leftmargin=2em]
    \item Select \(4\times4\) Sudoku or Boolean formula completion.
    \item Represent each variable as a finite categorical variable.
    \item Encode constraints as tensor factors.
    \item Condition on observed variables.
    \item Contract to obtain marginals.
    \item Iteratively fill high-confidence variables.
\end{enumerate}

\paragraph{Deliverables.}
\begin{itemize}[leftmargin=2em]
    \item script \texttt{tn\_masked\_completion.py};
    \item plots of violation \(V(x^{(k)})\);
    \item success rate vs masking ratio;
    \item contraction cost vs problem size.
\end{itemize}

\subsection{Milestone 2: Add noisy learned denoiser}

\paragraph{Objective.}
Test product-of-experts correction.

\paragraph{Steps.}
\begin{enumerate}[leftmargin=2em]
    \item Train a small MLP/Transformer denoiser on corrupted symbolic instances.
    \item At inference, compute both neural logits and TN marginals.
    \item Combine via
    \[
        \log \tilde p_i(v)
        =
        \beta \log p_\theta(v)
        +
        (1-\beta)\log q_i(v).
    \]
    \item Sweep \(\beta\).
\end{enumerate}

\paragraph{Expected outcome.}
Tensor-network marginals should improve constraint satisfaction, especially out of distribution or under high masking.

\subsection{Milestone 3: tinyTT adaptive mask policy}

\paragraph{Objective.}
Represent the mask-selection policy in TT format.

\paragraph{Steps.}
\begin{enumerate}[leftmargin=2em]
    \item Generate examples of partial states and oracle repair masks.
    \item Define candidate masks from local heuristics and random perturbations.
    \item Implement TT scoring \(S_\eta(m;x)\).
    \item Train with candidate-set cross entropy or pairwise ranking.
    \item Compare against local entropy remasking.
\end{enumerate}

\paragraph{Success criterion.}
The TT policy should reduce the number of denoising steps or improve final satisfaction at fixed denoising budget.

\subsection{Milestone 4: Bridge to a real masked diffusion LM}

\paragraph{Objective.}
Use an open MDLM checkpoint if available, or emulate one with a small masked Transformer.

\paragraph{Steps.}
\begin{enumerate}[leftmargin=2em]
    \item Choose a constrained text/code task.
    \item Extract structured variables from generated text.
    \item Use a verifier to identify defective spans.
    \item Remask defective spans.
    \item Run the MDLM repair step.
    \item Compare standard confidence unmasking against verifier-guided remasking.
\end{enumerate}

\paragraph{Practical starting point.}
Do not begin with open-ended natural language. Start with constrained domains:
\begin{itemize}[leftmargin=2em]
    \item arithmetic expressions;
    \item JSON/YAML repair;
    \item simple Python function repair with unit tests;
    \item theorem skeletons with symbolic placeholders;
    \item typed lambda-calculus expressions.
\end{itemize}

\section{Possible theoretical results}

\subsection{Residual-decrease theorem}

Formalize verifier-guided remasking as adaptive defect correction.

Let \(V(x)\) be a constraint violation functional. Let \(M(x)\) be a mask selector. Suppose the denoiser satisfies a local expected improvement estimate:
\[
    \E[V(D_\theta(x,M(x)))]
    \le
    V(x) - \alpha \Delta_M(x) + \varepsilon_\theta,
\]
where \(\Delta_M(x)\) is the removable violation under mask \(M\). If the mask selector is quasi-optimal,
\[
    \Delta_{M(x)}(x)
    \ge
    \gamma \max_{M\in\mathcal A}\Delta_M(x),
\]
and if the maximum removable violation controls the total violation,
\[
    \max_{M\in\mathcal A}\Delta_M(x)
    \ge
    c V(x),
\]
then
\[
    \E[V(x^{(k+1)})]
    \le
    (1-\alpha\gamma c)\E[V(x^{(k)})] + \varepsilon_\theta.
\]
This yields linear convergence up to model error:
\[
    \E[V(x^{(K)})]
    \le
    (1-\alpha\gamma c)^K V(x^{(0)})
    +
    \frac{\varepsilon_\theta}{\alpha\gamma c}.
\]

\subsection{Tensor-rank approximation of mask policies}

Let \(S^\star(m)\) be the optimal mask-value function. If \(S^\star\) has low effective TT rank \(r\), then there exists a TT policy with storage
\[
    \mathcal O(n r^2)
\]
for binary masks. Approximation error in the value function implies regret bounds for candidate-set selection:
\[
    S^\star(m^\star)-S^\star(\hat m)
    \le
    2\|S^\star-S_\eta\|_{\infty,\mathcal C}.
\]
This gives a clean reason to study low-rank mask policies.

\subsection{Tensor-network logical marginals as denoising priors}

Let \(q_i(v)\) be the exact logical marginal under constraints and observations. If the neural denoiser satisfies
\[
    \KL(q_i \| p_{\theta,i}) \le \delta_i,
\]
then product-of-experts correction can reduce logical inconsistency when \(q_i\) is sharper than \(p_{\theta,i}\). One possible route is to bound expected violation under
\[
    \tilde p_i(v) \propto p_{\theta,i}(v)^\beta q_i(v)^{1-\beta}.
\]
The analysis should exploit the fact that \(\tilde p_i\) interpolates between neural plausibility and logical consistency.

\section{Most promising paper direction}

The most coherent research paper would be:

\begin{quote}
\emph{Verifier-Guided Tensor-Structured Masked Diffusion for Logical Sequence Reasoning}
\end{quote}

The paper would contain:
\begin{enumerate}[leftmargin=2em]
    \item a formal view of MDLM inference as iterative masked defect correction;
    \item tensor-network logical marginals as inference-time guidance;
    \item a TT mask policy for adaptive remasking;
    \item residual-decrease theory under abstract denoiser assumptions;
    \item experiments on symbolic reasoning and program repair.
\end{enumerate}

This is more distinctive than merely compressing Transformer layers with tensor trains. The central novelty is to use tensor methods to control the denoising process.

\section{Suggested order of work}

\paragraph{Week 1.}
Implement pure \texttt{tnreason} masked completion on a small finite logical task. Produce plots of violation decay.

\paragraph{Week 2.}
Add a noisy denoiser and product-of-experts correction with TN marginals.

\paragraph{Week 3.}
Implement a \texttt{tinyTT} mask scorer and train it on oracle repair masks.

\paragraph{Week 4.}
Write the abstract residual-decrease theorem and test whether the empirical curves match its assumptions.

\paragraph{Week 5--6.}
Move to a constrained language/code task: JSON repair, arithmetic expression repair, or Python function repair with unit tests.

\paragraph{Week 7--8.}
Integrate a small masked Transformer or available MDLM checkpoint and compare standard confidence unmasking with verifier-guided tensor remasking.

\section{Risk assessment}

\paragraph{Low risk.}
Pure tensor-network masked inference, product-of-experts correction, and symbolic verifier-guided remasking.

\paragraph{Medium risk.}
TT mask policies. The main risk is that local heuristics may already perform well on small tasks. The task must require nonlocal mask dependencies.

\paragraph{High risk.}
Direct integration with a large diffusion LLM. Engineering friction, checkpoint availability, and tokenizer-level alignment may dominate.

\paragraph{Avoid initially.}
Do not start by tensorizing the full Transformer. This is technically possible but unlikely to reveal the conceptual benefit of combining diffusion LMs and logical tensor networks.

\section{Conclusion}

The most promising connection between diffusion LLMs and tensor methods is not simple model compression. The stronger idea is to view masked diffusion as recursive adaptive inference over a high-dimensional symbolic tensor state. Logical tensor networks can provide differentiable or discrete constraint residuals; tensor-network contractions can provide logical marginals; tensor trains can parameterize mask policies; and recursive remasking can turn diffusion generation into a verifier-guided correction process.

A practical first step is to build a small symbolic masked-diffusion analogue using \texttt{tnreason}, then add noisy denoising and tensor-network correction. In parallel, \texttt{tinyTT} can be used to test whether adaptive mask policies have exploitable low-rank structure. If these small experiments succeed, the next step is to attach the same machinery to a real masked diffusion LM or a small masked Transformer.
\end{document}
